\documentclass[11pt,a4paper]{article}

\usepackage[margin=1in]{geometry}
\usepackage{amsmath,amssymb,amsthm,mathtools}
\usepackage{booktabs}
\usepackage{microtype}
\usepackage[dvipsnames]{xcolor}
\usepackage[colorlinks=true,linkcolor=MidnightBlue,citecolor=MidnightBlue,urlcolor=MidnightBlue]{hyperref}
\usepackage{enumitem}

\theoremstyle{plain}
\newtheorem{theorem}{Theorem}[section]
\newtheorem{proposition}[theorem]{Proposition}
\newtheorem{lemma}[theorem]{Lemma}
\newtheorem{corollary}[theorem]{Corollary}

\theoremstyle{definition}
\newtheorem{definition}[theorem]{Definition}
\newtheorem{remark}[theorem]{Remark}

\newcommand{\R}{\mathbb{R}}
\newcommand{\C}{\mathbb{C}}
\newcommand{\N}{\mathbb{N}}
\newcommand{\Rnn}{\mathbb{R}_{\ge 0}}
\DeclareMathOperator{\RePart}{Re}
\DeclareMathOperator{\ImPart}{Im}
\DeclareMathOperator{\srank}{separableRank}
\DeclareMathOperator{\finrank}{finrank}
\DeclareMathOperator{\spanop}{span}
\DeclareRobustCommand{\leanid}[1]{{\def\_{\textunderscore\allowbreak}\texttt{#1}}}
\emergencystretch=1.5em

\title{\textbf{Separable Rank of Radial Kernels in Two and Higher Dimensions:\\
Machine-Checked Proofs in Lean~4 and Isabelle/HOL}}

\author{GitHub Copilot\\[3pt]
\small Source archive (CC0~1.0): \url{https://github.com/OlliNiemitalo/kernel-rank-lean}}

\date{November 11, 2026}

\begin{document}
\maketitle

\begin{abstract}
A $d$-variate kernel $K : X_0 \times \dots \times X_{d-1} \to \mathbb{K}$ has \emph{separable rank} at most $N$ if it is a sum of $N$ coordinatewise products $\sum_{n=0}^{N-1} \prod_{j=0}^{d-1} f_{n,j}(x_j)$. On a $d$-dimensional Cartesian grid of support width $L$, a rank-$R$ kernel evaluates in $d R L$ one-dimensional taps per grid point rather than $L^d$. While a positive Gaussian $\exp(-\alpha\|x\|^2)$ has separable rank $1$, approximating flat-top or oscillatory isotropic kernels---such as defocus disks, Airy point-spread functions, circular low-pass filters, or band-limited wave propagators---with finite coordinatewise separable rank relies on signed or oscillatory profiles built from complex-rate Gaussians and exponential polynomials.

We give a complete machine-checked characterization in Lean~4 (with Mathlib) of the $d$-way separable rank of continuous circularly and spherically symmetric kernels $K(x) = \varphi(\|x\|^2)$ on $\R^d$ for all $d \ge 2$:
(1)~$\srank(K) < \infty$ over $\R$ if and only if $\varphi$ agrees on $[0,\infty)$ with a well-formed exponential-polynomial representation $\sum_i \RePart(P_i(s)e^{\lambda_i s})$ (nonzero polynomials $P_i$, distinct rates $\lambda_i$ with $\ImPart\lambda_i \ge 0$, and $P_i \in \R[X]$ when $\ImPart\lambda_i = 0$);
(2)~over $\C$, a continuous radial profile $\psi(\|x\|^2)$ has exact $d$-way separable rank $\sum_{\mu \in S} ((d-1)\deg P_\mu + 1)$ for nonzero polynomials $P_\mu$ at distinct rates;
(3)~over $\R$, a continuous radial profile $\varphi(\|x\|^2)$ has exact $d$-way separable rank $\sum_i c_d(P_i,\lambda_i)$ where each real-rate component ($\ImPart\lambda_i = 0$) costs $(d-1)\deg P_i + 1$ real $d$-way products (achieved by real Lagrange interpolation) and each nonreal conjugate-pair component ($\ImPart\lambda_i > 0$) costs $R_i = 2(d-1)\deg P_i + d$ real $d$-way products, achieved via an alternating $R_i$-point root-of-unity polynomial-product identity (giving $d$ rather than $2^{d-1}$ real products for a complex Gaussian when $\deg P_i = 0$) and proved optimal by Alder--Strassen degree peeling over $\R$; and
(4)~in $d = 2$, scalar extension $\R \hookrightarrow \C$ preserves rank unconditionally, whereas for $d \ge 3$ it can strictly lower rank (e.g.\ $\RePart(e^{i\|x\|^2})$ has complex rank $2$ and real rank $d$, and $\|x\|^2\cos(\|x\|^2)$ in $d = 3$ has complex rank $6$, real border rank at most $6$, and exact real rank $7$); moreover, in $d = 2$ the exact rank equals the degree of the minimal monic constant-coefficient half-line differential annihilator $p(D)\varphi_{\C} = 0$ ($p \in \C[X]$).

Necessity reduces to showing that a continuous one-parameter semigroup $M : [0,\infty) \to \mathcal{A}$ is an exponential $M(t) = \exp(tA)$. Alongside the $96$-line Banach-algebra integral-smoothing proof used in the main development, we formalize four finite-dimensional alternatives---one by finite differences in Lean~4 and three by Euler's infinitesimal step in Lean~4 and Isabelle/HOL using a homogeneous common-scale normalization $\eta = \varepsilon + \|M(\varepsilon) - 1\|$ that avoids any \emph{a priori} difference-quotient bound---and compare their lengths, library reliance, and generality. Finally, we connect the root-of-unity identity to two-dimensional harmonic (``beach-ball'') differential operators and evaluate concrete pass-count reductions for steerable derivative filter banks and three-dimensional isotropic propagators.
\end{abstract}

\section{Introduction}

Convolving a $d$-dimensional array with a dense $L \times \dots \times L$ spatial kernel requires $L^d$ multiply--accumulate operations per output sample. When the kernel $K : \R^d \to \mathbb{K}$ admits a rank-$R$ coordinatewise decomposition
\begin{equation}\label{eq:sop}
  K(x_0, \dots, x_{d-1}) = \sum_{n=0}^{R-1} \prod_{j=0}^{d-1} f_{n,j}(x_j),
\end{equation}
the $d$-dimensional convolution reduces to $R$ cascades of $d$ one-dimensional convolutions, costing $d R L$ operations per sample (or $O(d R)$ operations independent of $L$ when the one-dimensional factors admit recursive causal/anti-causal realizations). The least $R \in \N \cup \{\infty\}$ for which~\eqref{eq:sop} holds is the \emph{separable rank} of $K$: matrix rank when $d = 2$, and functional canonical polyadic (CP) tensor rank when $d \ge 3$.

For circularly or spherically symmetric (\emph{radial}) kernels in dimension $d \ge 2$,
\[
  K(x) = \varphi\!\left(\sum_{j=0}^{d-1} x_j^2\right), \qquad x = (x_0,\dots,x_{d-1}) \in \R^d,
\]
Maxwell's characterization of the normal distribution implies that a continuous nonvanishing radial kernel of separable rank $1$ has a Gaussian profile $\varphi(s) = \alpha e^{\lambda s}$. Real-rate Gaussians ($\lambda < 0$), however, are strictly monotone and convex in $s = \|x\|^2$; nonnegative combinations of them cannot exhibit flat passbands, sharp cutoffs, or zero crossings. In image processing, optics, and numerical wave propagation, O.~Niemitalo~\cite{niemitalo2019lowpass,niemitalo_kernel_optimizer} and related filter-design work~\cite{lampert2006anisotropic} approximate non-Gaussian radial profiles---such as circular defocus disks, Airy diffraction patterns $(J_1(r)/r)^2$, and ideal circular low-pass filters $\omega_c J_1(\omega_c r)/(2\pi r)$---by sums of \emph{complex-rate} Gaussians $\RePart(\beta e^{(-\alpha + i\omega)\|x\|^2})$, whose damped radial oscillations $\cos(\omega\|x\|^2 + \theta)$ synthesize sharp transitions and sidelobe rings while keeping finite separable rank.

This raises three mathematical questions:
\begin{enumerate}[leftmargin=2em,itemsep=2pt]
  \item \textbf{Classification ($d \ge 2$):} Must \emph{every} continuous radial profile $\varphi : [0,\infty) \to \R$ of finite $d$-way separable rank be an exponential polynomial in $s = \|x\|^2$?
  \item \textbf{Exact $d$-way rank over $\C$ and over $\R$:} Given an exponential-polynomial or multi-Gaussian profile with distinct rates, what is the exact minimum number of $d$-way product terms in~\eqref{eq:sop}? Can algebraic recombinations beat naive real-and-imaginary expansion, and does complexification preserve separable rank beyond two dimensions?
  \item \textbf{Proof architecture in interactive theorem provers:} In formalizing necessity---whose analytic core is the continuous one-parameter semigroup theorem $M(t) = \exp(tA)$ on $[0,\infty)$---can an elementary infinitesimal proof in the style of Euler be made simpler than classical integral smoothing once formalized in Lean~4 or Isabelle/HOL?
\end{enumerate}

\subsection{Contributions and overview}

All numbered theorems, propositions, and lemmas in Sections~\ref{sec:core}--\ref{sec:semigroup} are machine-checked without unproved steps (\leanid{sorry} in Lean~4, oracles in Isabelle/HOL) in the open-source repository \texttt{kernel-rank-lean}, using Lean \texttt{v4.36.0-rc2} with Mathlib \texttt{v4.36.0-rc2} and Isabelle2026; Section~\ref{sec:applications} (along with the border-rank remark in Section~\ref{sec:nd}) develops analytical corollaries for harmonic differential operators and continuous-domain pass schedules:

\begin{itemize}[leftmargin=1.6em,itemsep=3pt]
  \item \textbf{Unified $d$-dimensional classification (Theorem~\ref{thm:classification}, \leanid{separableRank\_radial\_lt\_top\_iff\_exists}):} For every $d \ge 2$ and every $\varphi : \R \to \R$ continuous on $[0,\infty)$, the radial kernel $x \mapsto \varphi(\sum_{j=0}^{d-1} x_j^2)$ has finite $d$-way separable rank if and only if $\varphi$ agrees on $[0,\infty)$ with a canonical real exponential-polynomial representation $\sum_i \RePart(P_i(s)e^{\lambda_i s})$.
  \item \textbf{Exact $d$-D complex radial rank (Theorem~\ref{thm:complex_nd}, \leanid{separableRank\_radial\_complex\_eq\_iff}):} For $d \ge 2$ and $\psi : \R \to \C$ continuous on $[0,\infty)$, $\srank(x \mapsto \psi(\|x\|^2)) = N$ if and only if $\psi(s) = \sum_{\mu \in S} P_\mu(s)e^{\mu s}$ on $[0,\infty)$ for distinct rates $\mu \in S$ and nonzero polynomials $P_\mu$ satisfying $\sum_{\mu \in S} ((d-1)\deg P_\mu + 1) = N$. The upper bound uses shifted root-of-unity Fourier extraction on truncated exponential polynomials; the lower bound uses multi-mode Alder--Strassen degree peeling.
  \item \textbf{Exact $d$-D real radial rank for all degrees (Theorem~\ref{thm:real_nd}, \leanid{separableRank\_radial\_eq\_iff}):} For every $d \ge 2$, $\srank(x \mapsto \varphi(\|x\|^2)) = N$ over $\R$ if and only if $\varphi$ agrees on $[0,\infty)$ with a well-formed representation whose $d$-dimensional cost (\leanid{multiTotalCost}) equals $N$, counting $(d-1)\deg P_i + 1$ for real rates ($\ImPart\lambda_i = 0$, via real Lagrange interpolation) and $R_i = 2(d-1)\deg P_i + d$ for nonreal rates ($\ImPart\lambda_i > 0$, via an alternating root-of-unity polynomial-product identity \leanid{re\_coeff\_prod\_eq\_sum\_re} and Alder--Strassen peeling over $\R$). Specializing to $\deg P_i = 0$ gives $r + d c$ ($d$ real terms rather than $2^{d-1}$ per complex Gaussian, \leanid{separableRank\_radial\_gaussian\_profile}), and specializing to $d = 2$ recovers $\mathrm{totalCost}$ (\leanid{separableRank\_radial\_two\_eq\_iff}) alongside the minimal differential-annihilator characterization (Theorem~\ref{thm:annihilator}).
  \item \textbf{Comparative formalization of the semigroup step (Section~\ref{sec:semigroup}):} We compare the $96$-line Banach-algebra integral-smoothing proof against a $256$-line finite-difference proof in Lean~4 and three finite-dimensional nonstandard-analysis proofs ($206$ lines in Lean with ultrafilter sequences, $399$ lines in Lean with an ultrapower interface, and $275$ lines in Isabelle/HOL), all simplified via a homogeneous common-scale normalization $\eta = \varepsilon + \|M(\varepsilon) - 1\|$.
  \item \textbf{Harmonic differential operators, worked examples, and pass-count bounds (Section~\ref{sec:applications}):} Specializing the root-of-unity identity to equal directional-derivative factors recovers and unifies odd and even planar ``beach-ball'' harmonic operators $\RePart((\partial_x + i\partial_y)^m)$~\cite{niemitalo2019beachball,perona1995deformable}, yields exact continuous-domain 2-way separable ranks for analytically differentiated multi-Gaussian steerable filters~\cite{niemitalo2019angle}, establishes continuous-domain pass-count lower bounds in $2$-D alongside divide-and-conquer pass-tree upper bounds in $d$-D, and compares operation counts and numerical SVD/CP ranks across concrete 2-D and 3-D filter designs.
\end{itemize}

\subsection{Related work}

\paragraph{Radial characterizations and Gaussian sum approximations.}
Maxwell's kinetic-theory argument (1860) and Schoenberg's characterization of completely monotone radial functions~\cite{schoenberg1938metric} identify single Gaussians and positive Gaussian mixtures among radial kernels. For oscillatory or compactly supported radial kernels in scientific computing, Beylkin and Monz\'on~\cite{beylkin2005exponentials} approximate profiles and Green's functions by short sums of complex exponentials and Gaussians; in image filtering, Lampert and Wirjadi~\cite{lampert2006anisotropic} study anisotropic Gaussian separations, while Niemitalo~\cite{niemitalo2019lowpass,niemitalo_kernel_optimizer} fits circular defocus and low-pass profiles directly with complex-variance Gaussians. None of these works establishes a machine-checked converse for arbitrary signed finite-rank radial kernels or exact $d$-way real CP ranks across higher-degree components.

\paragraph{Hankel operators, Waring rank, and real vs.\ complex tensor rank.}
Reducing radial kernels via $s_j = x_j^2$ turns $K$ into the $d$-way additive Hankel kernel $\varphi(s_0 + \dots + s_{d-1})$. In two variables, Kronecker's theorem (1881) relates finite-rank discrete Hankel matrices to rational generating functions and linear recurrences; our continuous half-line version uses translation semigroups on $\mathcal{W}(\varphi)$ to avoid boundary cases and smoothness assumptions beyond continuity. In $d \ge 3$ variables, symmetric and Hankel tensor decompositions go back to Sylvester's apolarity theory (1851) and Comon and Mourrain~\cite{comon1996decomposition}; Landsberg~\cite{landsberg2012tensors} treats multilinear rank substitution (originating in Alder and Strassen~\cite{alder1981alg}) and the gap between border rank and exact rank, while Reznick~\cite{reznick2013length} and Qi, Comon, and Lim~\cite{qi2016semialgebraic} show that real symmetric/CP rank can strictly exceed complex rank. Our Theorems~\ref{thm:complex_nd} and~\ref{thm:real_nd} give exact closed-form complex and real $d$-way ranks for the entire continuous Hankel exponential-polynomial class, pinning down the exact real-over-complex gap $(\mathrm{cost}_{d,\R} - \mathrm{cost}_{d,\C} = d - 2$ for each nonreal conjugate pair, independent of polynomial degree!).

\paragraph{Formalized analysis and nonstandard analysis.}
Banach-space calculus, Bochner integration, and normed-algebra exponentials in Lean's Mathlib~\cite{mathlib2020} provide the infrastructure for \leanid{Semigroup.lean}, while Isabelle/HOL's multivariate analysis library~\cite{holzl2013type} and Fleuriot's formalization of nonstandard analysis~\cite{fleuriot2000mechanizing} (\texttt{HOL-Nonstandard\_Analysis}) underpin \leanid{Semigroup\_NSA.thy}. Section~\ref{sec:semigroup} provides a controlled side-by-side comparison of standard and nonstandard formalizations of the same theorem across both systems.

\section{Definitions, Surjections, and the Two-Dimensional Core}\label{sec:core}

\subsection{Multi-way separable rank and exponential-polynomial representations}

Throughout the paper, $\varphi : \R \to \R$ denotes a real-valued radial profile and $\psi : \R \to \C$ denotes a complex-valued radial profile. Let $\iota$ be a finite index type and $(X_i)_{i \in \iota}$ a family of sets, and let $\mathbb{K}$ be a commutative semiring. In \leanid{Defs.lean}:

\begin{definition}[\leanid{IsSumOfProducts}, \leanid{separableRank}]
A kernel $K : (\prod_{i \in \iota} X_i) \to \mathbb{K}$ is a \emph{sum of $N$ products} if there exist univariate factors $f : \operatorname{Fin} N \to \prod_{i \in \iota} (X_i \to \mathbb{K})$ such that
\[
  K(x) = \sum_{n \in \operatorname{Fin} N} \prod_{i \in \iota} f_n(i, x_i) \qquad \text{for all } x \in \prod_{i \in \iota} X_i.
\]
Its \emph{separable rank} is $\srank(K) = \inf\{N \in \N : \operatorname{IsSumOfProducts}(K, N)\} \in \N \cup \{\infty\}$ (\leanid{ENat}, written \texttt{Nat}-with-top in Mathlib).
\end{definition}

Separable rank is subadditive under pointwise sums (\leanid{IsSumOfProducts.finsetSum}) and monotone under coordinatewise composition (\leanid{IsSumOfProducts.comp}). Crucially, coordinatewise \emph{surjections} preserve separable rank exactly (\leanid{separableRank\_comp\_surjective}): if $g_i : Y_i \to X_i$ has a right inverse for each $i \in \iota$, then
\begin{equation}\label{eq:surj}
  \srank\!\bigl(y \mapsto K(i \mapsto g_i(y_i))\bigr) = \srank(K).
\end{equation}
Since $x \mapsto x^2$ maps $\R$ onto $\Rnn = [0,\infty)$, applying~\eqref{eq:surj} to $\iota = \operatorname{Fin} d$ identifies the $d$-way separable rank of any radial kernel with that of the $d$-way additive Hankel kernel (\leanid{separableRank\_radial\_eq\_hankel} in \leanid{Translations.lean}):
\begin{equation}\label{eq:hankel}
  \srank\!\left(x \in \R^d \mapsto \varphi\!\left(\sum_{j=0}^{d-1} x_j^2\right)\right)
  = \srank\!\left(s \in \Rnn^d \mapsto \varphi\!\left(\sum_{j=0}^{d-1} s_j\right)\right).
\end{equation}
The same lemma immediately transfers every theorem below to axis-aligned ellipsoidal kernels $\varphi(\sum_j x_j^2 / a_j^2)$ ($a_j \ne 0$) and $\ell_p$ superellipsoidal kernels $\varphi(\sum_j |x_j|^p)$ ($p > 0$). Moreover, for $d \ge 2$, padding $(y_0, y_1) \in \R^2$ with $d - 2$ zeros (\leanid{padTwo}) yields the slice monotonicity bound (\leanid{separableRank\_radial\_two\_le}):
\begin{equation}\label{eq:slice}
  \srank\!\left(y \in \R^2 \mapsto \varphi(y_0^2 + y_1^2)\right)
  \le \srank\!\left(x \in \R^d \mapsto \varphi\!\left(\sum_{j=0}^{d-1} x_j^2\right)\right).
\end{equation}

To state real-valued classifications without redundant conjugate duplicates, a \emph{component} $c = (P, \lambda)$ (\leanid{Component}) pairs a polynomial $P \in \C[X]$ with a complex rate $\lambda \in \C$ and evaluates on $\R$ as $c.\mathrm{eval}(s) = \RePart(P(s)e^{\lambda s})$, with profile $\mathrm{profile}(\mathrm{cs}, s) = \sum_{c \in \mathrm{cs}} c.\mathrm{eval}(s)$.

\begin{definition}[\leanid{Component.WellFormed}, \leanid{IsRepresentation}, \leanid{multiTotalCost}]
A component $c = (P,\lambda)$ is \emph{well-formed} if $P \ne 0$, $\ImPart\lambda \ge 0$, and $\ImPart\lambda = 0 \implies P \in \R[X]$ (every coefficient of $P$ has zero imaginary part). A list $\mathrm{cs}$ is a \emph{representation} (\leanid{IsRepresentation}) if all its components are well-formed and their rates $\lambda$ are pairwise distinct. In dimension $d \ge 2$, the $d$-way rendering cost of $c$ (\leanid{Component.multiCost}) is
\[
  c.\mathrm{multiCost}(d) = \begin{cases}
    (d - 1)\deg P + 1 & \text{if } \ImPart\lambda = 0,\\
    2(d - 1)\deg P + d & \text{if } \ImPart\lambda > 0,
  \end{cases}
\]
with list sum $\mathrm{multiTotalCost}(d, \mathrm{cs}) = \sum_{c \in \mathrm{cs}} c.\mathrm{multiCost}(d)$. At $d = 2$, $c.\mathrm{multiCost}(2) = c.\mathrm{cost}$ ($\deg P + 1$ for real rates, $2(\deg P + 1)$ for nonreal rates, summed in \leanid{totalCost}). A representation is \emph{Gaussian} (\leanid{IsGaussianRepresentation}) if every $c \in \mathrm{cs}$ has $\deg P = 0$, where $\mathrm{multiTotalCost}(d, \mathrm{cs}) = \mathrm{gaussianTotalCost}(d, \mathrm{cs}) = r + d c$ for $r$ real and $c$ nonreal components.
\end{definition}

\begin{remark}[Bijection with conjugate-symmetric complex exponential polynomials]
Why do the conditions $\ImPart\lambda \ge 0$ and $(\ImPart\lambda = 0 \implies P \in \R[X])$ make real representations canonical? Because $\RePart(\overline{P}(s)e^{\bar\lambda s}) = \RePart(P(s)e^{\lambda s})$ (where $\overline{P}$ conjugates coefficients) and $\RePart(i a e^{\mu s}) = 0$ for $a, \mu \in \R$, unrestricted pairs $(P, \lambda)$ over-parameterize real profiles. Under \leanid{Component.WellFormed}, each real rate ($\ImPart\lambda = 0, P \in \R[X]$) expands in $\C$ to a single term $P(s)e^{\lambda s}$, whereas each upper-half-plane rate ($\ImPart\lambda > 0$) expands via $\RePart z = \frac{1}{2}(z + \bar z)$ to the conjugate pair $\frac{1}{2}P(s)e^{\lambda s} + \frac{1}{2}\overline{P}(s)e^{\bar\lambda s}$ at distinct rates $\lambda \ne \bar\lambda$ (\leanid{injective\_rates\_with\_conjugates}). Since exponential polynomials at distinct complex rates are $\C$-linearly independent on $(0,\infty)$ (\leanid{eq\_zero\_of\_sum\_expPoly\_eq\_zero'} in \leanid{ExpPoly.lean}), a continuous real profile determines its unordered set of well-formed components uniquely. Notice also that for $\ImPart\lambda > 0$, the complex rank of this conjugate pair by Theorem~\ref{thm:complex_nd} is $2((d - 1)\deg P + 1) = 2(d - 1)\deg P + 2$, so $c.\mathrm{multiCost}(d) - 2((d - 1)\deg P + 1) = d - 2$: every nonreal conjugate pair incurs an exact real-over-complex rank penalty of $d - 2$, independent of $\deg P$.
\end{remark}

\subsection{Two-coordinate row span, complexification, and basis-free necessity}

When $\iota = \operatorname{Fin} 2$ and $\mathbb{K}$ is a field, \leanid{RowSpan.lean} defines $\operatorname{rowSpan}(K) \le (X \to \mathbb{K})$ as the $\mathbb{K}$-linear span of the sections $\{x \mapsto K(x, y) : y \in X\}$ (fixing the second coordinate and varying the first) and proves that $\srank(K) < \infty$ holds if and only if $\operatorname{rowSpan}(K)$ is finite-dimensional over $\mathbb{K}$, in which case $\srank(K) = \finrank_{\mathbb{K}}(\operatorname{rowSpan} K)$ (\leanid{separableRank\_eq\_finrank\_rowSpan}). For the two-variable Hankel kernel $(s_0, s_1) \mapsto \varphi(s_0 + s_1)$ on $\Rnn^2$, the row span is literally the \emph{translate span}
\[
  \mathcal{W}(\varphi) = \spanop_{\mathbb{K}}\{u \in \Rnn \mapsto \varphi(u + t) : t \in \Rnn\} \le (\Rnn \to \mathbb{K}).
\]

Scalar extension $\R \hookrightarrow \C$ (\texttt{Complexification.lean}) gives the bound $\srank_{\C}(K_{\C}) \le \srank_{\R}(K)$ for every finite index set $\iota$, where $K_{\C}(x) = (K(x) : \C)$ (\leanid{separableRank\_ofReal\_le}). In \emph{two} coordinates ($\iota = \operatorname{Fin} 2$), the reverse inequality also holds (\leanid{separableRank\_ofReal}): $\R$-linearly independent rows in $\R^X$ remain $\C$-linearly independent in $\C^X$ (tested by taking real and imaginary parts of any $\C$-linear relation), so $\finrank_{\C}\operatorname{rowSpan}(K_{\C}) = \finrank_{\R}\operatorname{rowSpan}(K)$, including when both ranks are $\infty$.

Necessity and the two-dimensional lower bound share a single coordinate-free realization in \leanid{TranslationSemigroup.lean}:
\begin{enumerate}[leftmargin=2em,itemsep=2pt]
  \item Suppose $\psi : \R \to \C$ is continuous on $[0,\infty)$ and $\dim_{\C}\mathcal{W}(\psi) < \infty$. Equipping $\mathcal{W}(\psi)$ with a norm transferred across a linear equivalence $\mathcal{W}(\psi) \simeq_{\C} \C^n$ makes $\mathcal{W}(\psi)$ a finite-dimensional Banach space whose induced topology agrees with the subspace topology inherited from the product topology on $\Rnn \to \C$ (since all Hausdorff vector topologies on $\C^n$ coincide).
  \item Because $\mathcal{W}(\psi)$ is invariant under non-negative shifts $(T_t f)(u) = f(u + t)$ ($t \ge 0$), continuity of $\psi$ gives pointwise continuity of $t \mapsto T_t f$ on each generator and hence operator-norm continuity of the semigroup $t \mapsto T_t \in \mathcal{L}(\mathcal{W}(\psi))$ on $[0,\infty)$ (\leanid{continuous\_shiftCLM}).
  \item Applying the Banach-algebra semigroup theorem (\leanid{Semigroup.lean}, Section~\ref{sec:semigroup}) yields a generator $A \in \mathcal{L}(\mathcal{W}(\psi))$ with $T_t = \exp(tA)$ for all $t \ge 0$, so $\psi(t) = \operatorname{eval}_0(\exp(tA)\psi_0)$ where $\psi_0 = \psi|_{\Rnn} \in \mathcal{W}(\psi)$ (\leanid{profile\_of\_generator}).
  \item Decomposing $\psi_0$ into generalized eigenspaces $\ker((A - \mu I)^k)$ of $A$ (\texttt{Module.End.}\allowbreak\leanid{iSup\_maxGenEigenspace\_eq\_top} in \leanid{Necessity.lean}) truncates the exponential series of $(A - \mu I)$ on each summand, expressing $\psi(t) = \sum_{\mu \in S} Q_\mu(t)e^{\mu t}$ for $t \ge 0$. When $\psi = \varphi_{\C}$ comes from a real profile $\varphi$, averaging $\varphi(t) = \frac{1}{2}(\varphi(t) + \overline{\varphi(t)})$ and pairing $\mu$ with $\bar\mu$ using exponential-polynomial uniqueness on $(0,\infty)$ (\leanid{ExpPoly.lean}, \leanid{RealExpPoly.lean}) regroups the sum into a well-formed real representation (\leanid{Representation.lean}).
\end{enumerate}

For the two-dimensional exact rank and annihilator characterization (\leanid{Differential.lean}, \leanid{ComplexRank.lean}, \leanid{Annihilator.lean}), let $D$ denote differentiation within $[0,\infty)$ (right differentiation at $0$). Applying Cayley--Hamilton to the generator $A$ on $\mathcal{W}(\psi)$ shows that the monic characteristic polynomial $\chi_A(X)$ of degree $\dim_{\C}\mathcal{W}(\psi)$ annihilates every orbit $t \mapsto \ell(\exp(tA)v)$, so $\chi_A(D)\psi = 0$ on $[0,\infty)$. Conversely, for $\psi(t) = \sum_{i} P_i(t)e^{\lambda_i t}$ with distinct rates $\lambda_i$ and $P_i \ne 0$, Hasse-derivative independence shows that $q(D)\psi = 0$ on $[0,\infty)$ if and only if the product annihilator
\[
  p(X) = \prod_i (X - \lambda_i)^{\deg P_i + 1}
\]
divides $q(X)$ (\leanid{differentialOperator\_sum\_expPoly\_eq\_zero\_iff}). Thus $p \mid \chi_A$, giving $\sum_i (\deg P_i + 1) = \deg p \le \deg\chi_A = \dim_{\C}\mathcal{W}(\psi)$, while Taylor shift expansion of $P_i(u + t)e^{\lambda_i(u+t)}$ gives the matching upper bound $\dim_{\C}\mathcal{W}(\psi) \le \sum_i (\deg P_i + 1)$.

\section{Exact Multi-Way Separable Ranks in $d \ge 2$ Dimensions}\label{sec:nd}

In dimension $d \ge 3$, there is no single row-span dimension equal to $\srank(K)$, and scalar extension $\R \hookrightarrow \C$ no longer preserves rank. Instead, \leanid{MultiComplexRank.lean} and \leanid{MultiRealRank.lean} establish exact $d$-way upper bounds by root-of-unity and Lagrange identities and matching lower bounds by mapping functional decompositions onto finite-dimensional multilinear forms.

\subsection{Bridging function spaces and multilinear tensor rank}

Suppose $V$ is a $\mathbb{K}$-vector space over a field $\mathbb{K}$, $v : \alpha \to V$ is an evaluation map whose image spans $V$ ($\spanop_{\mathbb{K}}(\operatorname{range} v) = \top$), and $M : V^d \to \mathbb{K}$ is a $d$-multilinear map (\leanid{MultilinearMap} $\mathbb{K}$ $(\operatorname{Fin} d \to V)$ $\mathbb{K}$). Lemma~\leanid{exists\_linear\_factors\_of\_span\_eq\_top} proves that if the composed kernel $s \in \alpha^d \mapsto M(v(s_0), \dots, v(s_{d-1}))$ admits a functional decomposition~\eqref{eq:sop} of rank $R$ with arbitrary set-theoretic factors $f_{n,j} : \alpha \to \mathbb{K}$, then there exist \emph{linear} functionals $\ell_{n,j} \in V^*$ such that
\begin{equation}\label{eq:multilinear_decomp}
  M(u_0, \dots, u_{d-1}) = \sum_{n=0}^{R-1} \prod_{j=0}^{d-1} \ell_{n,j}(u_j) \qquad \text{for all } u_0, \dots, u_{d-1} \in V.
\end{equation}
Indeed, selecting a $\mathbb{K}$-basis $\{v(s_b)\}_{b \in B}$ of $V$ from $\operatorname{range}(v)$ defines unique linear functionals $\ell_{n,j} \in V^*$ satisfying $\ell_{n,j}(v(s_b)) = f_{n,j}(s_b)$; both sides of~\eqref{eq:multilinear_decomp} are $d$-multilinear maps that agree on every $d$-tuple of basis vectors, so they agree on $V^d$. To verify $\spanop_{\mathbb{K}}(\operatorname{range} v) = \top$ when $\dim_{\mathbb{K}} V < \infty$, Lemma~\leanid{span\_range\_eq\_top\_of\_dual} shows by quotient duality that it suffices to check $L \circ v = 0 \implies L = 0$ for every $L \in V^*$.

\subsection{Exact $d$-way complex exponential-polynomial rank}

Fix $d \ge 2$, a finite index set $\iota$, injective rates $\lambda : \iota \to \C$, and nonzero polynomials $P_i \in \C[X]$ of degree $k_i = \deg P_i$ ($i \in \iota$).

\paragraph{Upper bound via shifted root-of-unity interpolation.}
For a single polynomial $P \in \C[X]$ of degree $k = \deg P$, put $L = (d - 1)k + 1 \ge 1$ and $\zeta = \exp(2\pi i / L)$. Using the $q$-th Hasse derivative $D^{[q]} P$ (so $P(s_0 + t) = \sum_{q=0}^k (D^{[q]} P)(s_0)\, t^q$), define the reversed Hasse polynomial on coordinate $0$ and the truncated exponential polynomials on coordinates $j \in \{1,\dots,d-1\}$ (\leanid{revHassePoly}, \leanid{truncExpPoly}):
\begin{gather*}
  R(s_0; X) = \sum_{q=0}^k q!\,(D^{[q]} P)(s_0)\, X^{k - q}, \qquad
  E_k(s_j; X) = \sum_{q=0}^k \frac{s_j^q}{q!} X^q,\\
  Q(X) = R(s_0; X) \prod_{j=1}^{d-1} E_k(s_j; X).
\end{gather*}
Binomial convolution gives $[X^q] \prod_{j=1}^{d-1} E_k(s_j; X) = (s_1 + \dots + s_{d-1})^q / q!$ for all $0 \le q \le k$, so the coefficient of $X^k$ in $Q(X)$ is precisely $P(s_0 + s_1 + \dots + s_{d-1})$. Crucially, because $d \ge 2$, we have $k < L$ and $\deg Q \le k + (d - 1)k = k + L - 1 < k + L$. For \emph{any} polynomial $Q$ with $\deg Q < k + L$ and $k < L$, discrete Fourier averaging over only $L$ roots of unity extracts $[X^k] Q(X)$ (\leanid{coeff\_eq\_sum\_rootsOfUnity}):
\[
  [X^k] Q(X) = \frac{1}{L} \sum_{r=0}^{L-1} \zeta^{-r k}\, Q(\zeta^r),
\]
because every monomial index $n \in \{0,\dots,k+L-1\}$ satisfies $-L < -k \le n - k < L$, so $L \mid (n - k)$ holds if and only if $n = k$. Multiplying the $j$-th factor of $Q(\zeta^r)$ by $e^{\lambda s_j}$ (and absorbing $\frac{1}{L}\zeta^{-r k}$ into coordinate $0$) expresses $P(\sum_{j=0}^{d-1} s_j)e^{\lambda\sum_{j=0}^{d-1} s_j}$ as a sum of $L = (d - 1)\deg P + 1$ coordinatewise products (\leanid{isSumOfProducts\_expPoly}).

\paragraph{Lower bound via Alder--Strassen degree peeling.}
Let $m_0(i) = \deg P_i + 1$ and form the coefficient evaluation space $V = \prod_{i \in \iota} \C^{m_0(i)}$ with $v(s)_{i,q} = \frac{s^q}{q!}e^{\lambda_i s}$ ($0 \le q < m_0(i)$). By univariate exponential-polynomial uniqueness (\leanid{eq\_zero\_of\_sum\_expPoly\_eq\_zero'}), no nonzero functional in $V^*$ vanishes on $\operatorname{range}(v)$, so $\spanop_{\C}(\operatorname{range} v) = \top$ (\leanid{span\_expMonomial\_eq\_top}). Under $v$, the Hankel kernel corresponds to the $d$-multilinear truncated-convolution form (\leanid{truncForm}) with cutoffs $c_i \le m_0(i)$,
\[
  \mathcal{T}_c(u^{(0)}, \dots, u^{(d-1)})
  = \sum_{i \in \iota}\; \sum_{q_0 + \dots + q_{d-1} < c_i} w_i(q_0 + \dots + q_{d-1}) \prod_{j=0}^{d-1} u^{(j)}_{i, q_j},
\]
where $w_i(c_i - 1) \ne 0$ whenever $c_i > 0$. Lemma~\leanid{truncForm\_rank\_ge} proves
\[
  \sum_{i :\, c_i > 0} \bigl((d - 1)(c_i - 1) + 1\bigr) \le R
\]
for every rank-$R$ decomposition~\eqref{eq:multilinear_decomp} by strong induction on $S = \sum_{i \in \iota} c_i$:
\begin{itemize}[leftmargin=1.6em,itemsep=2pt]
  \item \emph{Case $c_{i_0} = 1$:} Evaluating $\mathcal{T}_c$ at $(e_{i_0,0}, \dots, e_{i_0,0})$ gives $w_{i_0}(0) \ne 0$, selecting an index $n_0$ with $\ell_{n_0,0}(e_{i_0,0}) \ne 0$. Let $\pi_{\ne i_0} : V \to V$ zero out block $i_0$. Substituting $\varphi_0(x) = \pi_{\ne i_0}(x) - \frac{\ell_{n_0,0}(\pi_{\ne i_0}(x))}{\ell_{n_0,0}(e_{i_0,0})} e_{i_0,0}$ on mode $0$ and $\pi_{\ne i_0}$ on modes $j \ge 1$ annihilates $\ell_{n_0,0} \circ \varphi_0$ while reducing $c_{i_0}$ from $1$ to $0$ (since $d \ge 2$ zeroes block $i_0$ on mode $1$), leaving at most $R - 1$ terms.
  \item \emph{Case $c_{i_0} \ge 2$:} Set $s_{\mathrm{top}} = e_{i_0, c_{i_0}-1}$ and $s_0 = e_{i_0, 0}$ ($s_{\mathrm{top}} \ne s_0$). Because $(c_{i_0} - 1) + (c_{i_0} - 1) \ge c_{i_0}$, any rank-$1$ perturbation of the form $\psi(x) = x - \gamma(x) s_{\mathrm{top}}$ leaves $\mathcal{T}_c$ unchanged whenever another coordinate equals $s_{\mathrm{top}}$. Stepping through modes $j = 0, \dots, d - 2$ sequentially evaluates $\mathcal{T}_c$ with $u^{(j)} = s_{\mathrm{top}}$ and $u^{(k)} = \psi_k(s_0)$ ($k \ne j$) to find $d - 1$ distinct indices $D \subseteq \operatorname{Fin} R$ ($|D| = d - 1$) killed on modes $0, \dots, d - 2$ by rank-$1$ perturbations $\psi_j$ along $s_{\mathrm{top}}$. On modes $0, \dots, d - 2$ we then truncate block $i_0$ to indices $q < c_{i_0} - 1$ before applying $\psi_j$, while on mode $d - 1$ we shift block $i_0$'s degree indices up by $+1$ ($e_{i_0, q} \mapsto e_{i_0, q+1}$ for $q < c_{i_0} - 1$): this transforms $\mathcal{T}_c$ into the reduced form $\mathcal{T}_{c'}$ ($c' = c[i_0 \mapsto c_{i_0} - 1]$) with shifted weights $w'_{i_0}(m) = w_{i_0}(m + 1)$---preserving the nonzero top weight $w'_{i_0}(c'_{i_0} - 1) = w_{i_0}(c_{i_0} - 1) \ne 0$---while dropping all $d - 1$ terms in $D$.
\end{itemize}
Combining this lower bound at $c = m_0$ with the upper bound and Section~\ref{sec:core} gives:

\begin{theorem}[$d$-D Complex Radial Rank, \leanid{separableRank\_radial\_complex\_eq\_iff}]\label{thm:complex_nd}
Let $d \ge 2$ and let $\psi : \R \to \C$ be continuous on $[0,\infty)$. For every $N \in \N$,
\[
  \srank\!\left(x \in \R^d \mapsto \psi\!\left(\sum_{j=0}^{d-1} x_j^2\right)\right) = N
\]
if and only if there exist a finite set $S \subset \C$ and polynomials $(P_\mu)_{\mu \in S}$ with $P_\mu \ne 0$ such that $\psi(s) = \sum_{\mu \in S} P_\mu(s)e^{\mu s}$ for all $s \ge 0$ and $\sum_{\mu \in S} \bigl((d - 1)\deg P_\mu + 1\bigr) = N$.
\end{theorem}

\subsection{Exact $d$-way real rank for arbitrary degrees (\leanid{MultiRealRank.lean})}

Before treating polynomial degrees, note that for constant factors ($d \ge 1$), expanding $\RePart(\prod_{j=0}^{d-1} z_j)$ into real and imaginary parts gives $2^{d-1}$ real products, whereas $d$ real products always suffice:

\begin{proposition}[Alternating $2d$-th root-of-unity identity, \leanid{re\_prod\_eq\_sum\_re}]\label{prop:roots}
For every integer $d \ge 1$ and all $z_0, \dots, z_{d-1} \in \C$,
\begin{equation}\label{eq:re_prod}
  \RePart\!\left(\prod_{j=0}^{d-1} z_j\right)
  = \frac{2^{d-1}}{d} \sum_{m=0}^{d-1} (-1)^m \prod_{j=0}^{d-1} \RePart\!\left(z_j e^{-i\pi m / d}\right).
\end{equation}
\end{proposition}

\begin{proof}
Write $\RePart(z_j e^{-i\pi m/d}) = \frac{1}{2}\bigl(\overline{z_j} e^{i\pi m/d} + z_j e^{-i\pi m/d}\bigr)$ and expand over subsets $S \subseteq \{0,\dots,d-1\}$: multiplying by $(-1)^m = e^{i\pi m}$ yields the geometric progression $w_S^m$ with $w_S = e^{2\pi i |S|/d}$ ($w_S^d = 1$), whose sum over $0 \le m < d$ equals $d$ for $S \in \{\emptyset, \{0,\dots,d-1\}\}$ and $0$ for $0 < |S| < d$.
\end{proof}

Absorbing $\frac{2^{d-1}}{d}(-1)^m$ into coordinate $0$ shows that the real part of a rank-$N$ complex $d$-way kernel has real rank at most $N d$ (\leanid{IsSumOfProducts.re}), establishing qualitative $d$-D classification:

\begin{theorem}[$d$-D Finite-Rank Classification, \leanid{separableRank\_radial\_lt\_top\_iff\_exists}]\label{thm:classification}
Let $d \ge 2$ and let $\varphi : \R \to \R$ be continuous on $[0,\infty)$. Then
\begin{align*}
  &\srank\!\left(x \in \R^d \mapsto \varphi\!\left(\sum_{j=0}^{d-1} x_j^2\right)\right) < \infty\\
  &\qquad\iff \exists\,\mathrm{cs},\ \operatorname{IsRepresentation}(\mathrm{cs}) \wedge \forall s \ge 0,\ \varphi(s) = \mathrm{profile}(\mathrm{cs}, s).
\end{align*}
\end{theorem}

For degree $k = \deg P \ge 1$, however, multiplying Theorem~\ref{thm:complex_nd}'s complex rank $(d-1)k + 1$ by $d$ via \leanid{IsSumOfProducts.re} gives $d(d-1)k + d$, which is non-optimal. Instead, sharp upper bounds differ for real and nonreal rates (\leanid{isSumOfProducts\_radial\_multiComponent}):
\begin{itemize}[leftmargin=1.6em,itemsep=2pt]
  \item \emph{Real rates ($\ImPart\lambda = 0$, $P \in \R[X]$ of degree $k$):} On coordinates $1, \dots, d-1$, the real polynomial $A(X) = \prod_{j=1}^{d-1} E_k(x_j^2; X) \in \R[X]$ has $\deg A \le (d - 1)k < L := (d - 1)k + 1$. Real Lagrange interpolation (\leanid{Lagrange.eq\_interpolate}) at the $L$ real nodes $0, 1, \dots, L - 1 \in \R$ expresses each $[X^q]A(X) = (\sum_{j=1}^{d-1} x_j^2)^q / q!$ ($0 \le q \le k$) as an $\R$-linear combination of $\prod_{j=1}^{d-1} E_k(x_j^2; r)$, yielding $(d - 1)k + 1$ real $d$-way products.
  \item \emph{Nonreal rates ($\ImPart\lambda > 0$, $P \in \C[X]$ of degree $k$):} We generalize Proposition~\ref{prop:roots} from scalars to polynomials of degree at most $k$ ($d \ge 2$):
\end{itemize}

\begin{proposition}[Universal root-of-unity polynomial-product identity, \leanid{re\_coeff\_prod\_eq\_sum\_re}]\label{prop:poly_roots}
Let $d \ge 2$, $k \ge 0$, and $R = 2(d - 1)k + d$. For any $U_0, \dots, U_{d-1} \in \C[X]$ with $\deg U_j \le k$,
\begin{equation}\label{eq:re_coeff_prod}
  \RePart\!\left([X^k]\prod_{j=0}^{d-1} U_j(X)\right)
  = \frac{2^{d-1}}{R} \sum_{r=0}^{R-1} (-1)^r \prod_{j=0}^{d-1} \RePart\!\left(U_j\!\left(e^{\frac{2\pi i r}{R}}\right) e^{-\frac{i\pi r (2k + 1)}{R}}\right).
\end{equation}
\end{proposition}

\begin{proof}
Expanding $\RePart w = \frac{1}{2}(\bar w + w)$ over subsets $S \subseteq \{0,\dots,d-1\}$ of size $s = |S|$ and coefficient indices $q \in \{0,\dots,k\}^d$, and noting that $(-1)^r = e^{i\pi r R / R}$ with $R = (2k + 1)d - 2k$, gives the $r$-th phase $\exp(2\pi i r K(S,q) / R)$ where
\[
  K(S, q) = (2k + 1)s + \sum_{j \notin S} q_j - \sum_{j \in S} q_j - k \in \mathbb{Z}.
\]
Because $0 \le q_j \le k$ and $d \ge 2$: (i)~when $S = \emptyset$ ($s = 0$), $-R < -k \le K(\emptyset, q) = \sum_j q_j - k \le (d-1)k < R$, so $R \mid K(\emptyset, q) \iff \sum_j q_j = k$; (ii)~when $S = \{0,\dots,d-1\}$ ($s = d$), $K = R + (k - \sum_j q_j)$, so $R \mid K \iff \sum_j q_j = k$; and (iii)~when $1 \le s \le d - 1$, $K(S,q) \ge (2k+1)s - s k - k = (k+1)s - k \ge 1$ and $K(S,q) \le (2k+1)s + (d-s)k - k = (k+1)s + (d-1)k \le 2(d-1)k + d - 1 = R - 1$, so $R \nmid K(S,q)$ and the geometric sum over $0 \le r < R$ vanishes identically. Applying~\eqref{eq:re_coeff_prod} to $U_0(X) = e^{\lambda x_0^2} R(x_0^2; X)$ and $U_j(X) = e^{\lambda x_j^2} E_k(x_j^2; X)$ ($j \ge 1$) yields $2(d - 1)k + d$ real $d$-way products.
\end{proof}

\paragraph{Unified Alder--Strassen lower bound over $\R$ (\leanid{realExpForm\_rank\_ge}).}
Let $G = (\prod_{i \in \iota_{\R}} \R^{m_{\R}(i)}) \times (\prod_{k \in \iota_{\C}} \C^{m_{\C}(k)})$ with $m(i) = \deg P_i + 1$. Conjugate-rate uniqueness shows via duality that the exponential-monomial curve $s \mapsto ((s^q/q!\,e^{\mu_i s}), (s^q/q!\,e^{\lambda_k s}))$ spans $G$ over $\R$ (\leanid{span\_realExpEval\_eq\_top}), transferring any rank-$R$ decomposition of the Hankel kernel to $\R$-linear functionals $\ell_{n,j} \in G^*$ for the combined truncated-convolution form \leanid{realExpForm}. Note first that every $L \in \operatorname{Hom}_{\R}(\C, \R)$ has a nonzero kernel vector $w \in \ker L \setminus \{0\}$ (\leanid{exists\_ker\_ne\_zero}), from which induction on $d \ge 2$ gives $d \le R$ for any single nonzero complex product $\RePart(\beta \prod_{j=0}^{d-1} z_j)$ (\leanid{single\_complex\_rank\_ge}). Strong induction on the cutoff sum then peels:
\begin{enumerate}[leftmargin=2em,itemsep=2pt]
  \item \emph{Real cutoffs (\leanid{realExpForm\_rank\_ge}):} Identical to Section~3.2 ($1$ term when $c_{\R}(i_0) = 1$, $d - 1$ terms when $c_{\R}(i_0) \ge 2$), while keeping the $\C$-coordinates of $G$ fixed so all nonreal blocks are untouched.
  \item \emph{Nonreal cutoff $c_{\C}(k_0) = 1$ (\leanid{cplxTruncForm\_rank\_ge}):} On mode $0$, evaluating at $1_{k_0}$ and at $w_1 \in \ker(\ell_{n_1,0}|_{\C e_{k_0,0}}) \setminus \{0\}$ kills $2$ distinct terms, and on each mode $j = 1, \dots, d - 2$, \leanid{single\_complex\_rank\_ge} forces one additional active term outside the killed set $D$ (else $|D| = j + 1 < d$ terms would represent $\RePart(\beta_{k_0}\prod z_q)$), reaching $|D| = d$ before zeroing block $k_0$ on mode $d - 1$.
  \item \emph{Nonreal cutoff $c_{\C}(k_0) \ge 2$:} Let $q_{\mathrm{top}} = c_{\C}(k_0) - 1 \ge 1$ with weight $\beta_{\mathrm{top}} \ne 0$. In each mode $j = 0, \dots, d - 2$, evaluating at $u(z)$ with $\operatorname{emb}_{\mathrm{top}}(z)$ on mode $j$, $\operatorname{emb}_0(\overline{\beta_{\mathrm{top}} z})$ on mode $d - 1$, and $\psi_m(\operatorname{emb}_0(1))$ elsewhere gives $|\beta_{\mathrm{top}} z|^2 > 0$ for every $z \in \C \setminus \{0\}$ (since $2 q_{\mathrm{top}} \ge c_{\C}(k_0)$ kills all top-degree cross-terms). Testing $z = 1$ and $z = w_1 \in \ker(\ell_{n_1,j} \circ \operatorname{emb}_{\mathrm{top}}) \setminus \{0\}$ kills $2$ distinct terms per mode---$2(d - 1)$ terms across modes $0, \dots, d - 2$---after which low-degree truncation on modes $0, \dots, d - 2$ and $+1$ index shifting on mode $d - 1$ reduce $c_{\C}(k_0) \mapsto c_{\C}(k_0) - 1$ with shifted top weight $\beta_{\mathrm{top}} \ne 0$.
\end{enumerate}

\begin{theorem}[Exact $d$-D Real Radial Rank, \leanid{separableRank\_radial\_eq\_iff}]\label{thm:real_nd}
Let $d \ge 2$ and let $\varphi : \R \to \R$ be continuous on $[0,\infty)$. For every $N \in \N$,
\begin{align*}
  &\srank\!\left(x \in \R^d \mapsto \varphi\!\left(\sum_{j=0}^{d-1} x_j^2\right)\right) = N\\
  &\iff \exists\,\mathrm{cs},\ \operatorname{IsRepresentation}(\mathrm{cs})
    \wedge (\forall s \ge 0,\ \varphi(s) = \mathrm{profile}(\mathrm{cs}, s))
    \wedge \mathrm{multiTotalCost}(d, \mathrm{cs}) = N.
\end{align*}
In particular, for a Gaussian representation with $r$ real and $c$ nonreal rates, $\srank = r + d c$ (\leanid{separableRank\_radial\_gaussian\_profile}), and in $d = 2$, $\srank = \mathrm{totalCost}(\mathrm{cs})$ (\leanid{separableRank\_radial\_two\_eq\_iff}).
\end{theorem}

\begin{theorem}[Minimal Differential Annihilator]\label{thm:annihilator}
Let $\varphi : \R \to \R$ be continuous on $[0,\infty)$ with finite-dimensional translate span and two-dimensional radial rank $\srank(x \in \R^2 \mapsto \varphi(x_0^2 + x_1^2)) = N < \infty$ (\leanid{exists\_minimal\_differential\_annihilator}). There exists a monic polynomial $p \in \C[X]$ of degree $N$ such that for every $q \in \C[X]$, the half-line differential equation $q(D)\varphi_{\C} = 0$ holds on $[0,\infty)$ (where $\varphi_{\C}(s) = (\varphi(s) : \C)$) if and only if $p \mid q$.
\end{theorem}

\paragraph{Worked example ($d = 3, \deg P = 1$): exact rank vs.\ complex rank and border rank.}
Consider $K(x_0, x_1, x_2) = (x_0^2 + x_1^2 + x_2^2)\cos(x_0^2 + x_1^2 + x_2^2)$ on $\R^3$, a single well-formed component $P(s) = s$, $\lambda = i$ ($\deg P = 1, \ImPart\lambda = 1 > 0$):
\begin{itemize}[leftmargin=1.6em,itemsep=2pt]
  \item \emph{Complex rank $= 6$:} Over $\C$, $K_{\C}$ splits into $\frac{1}{2}s e^{is} + \frac{1}{2}s e^{-is}$, giving exact complex 3-way rank $((3-1)\cdot 1 + 1) + ((3-1)\cdot 1 + 1) = 6$ by Theorem~\ref{thm:complex_nd}.
  \item \emph{Pointwise limit of rank-$6$ kernels (border rank $\le 6$):} As an unformalized analytical contrast, for every $h > 0$ the difference quotient $\frac{1}{h}\bigl[\RePart\prod_{j=0}^2 (1 + h x_j^2)e^{i x_j^2} - \RePart\prod_{j=0}^2 e^{i x_j^2}\bigr]$ converges pointwise (and uniformly on compact sets) to $K(x)$ as $h \to 0$, and by Proposition~\ref{prop:roots} each $\RePart(\prod_{j=0}^2 \cdots)$ uses $3$ real 3-way products, totaling $3 + 3 = 6$ real products at every $h > 0$ with diverging $O(h^{-1})$ weights.
  \item \emph{Exact real rank $= 7$:} By Theorem~\ref{thm:real_nd}, $\srank_{\R}(K) = 2(3 - 1)\cdot 1 + 3 = 7$ (achieved with bounded weights by Proposition~\ref{prop:poly_roots} at $R = 7$, vs.\ $9$ from naive \leanid{IsSumOfProducts.re}). To see why $6$ terms cannot represent $K$ exactly, pass to $W = \C[T]/(T^2) \cong \C \oplus \C\varepsilon$ ($\varepsilon^2 = 0$) where $\mathcal{M}(u, v, w) = \RePart(u_1 v_0 w_0 + u_0 v_1 w_0 + u_0 v_0 w_1)$. On $E = 0 \oplus \C\varepsilon \cong \R^2$, the slice $\mathcal{M}((0, z_1), (v_0, 0), (w_0, 0)) = \RePart(z_1 v_0 w_0)$ forces the mode-$0$ functionals $\{\ell_n|_E\}$ to span $E^*$ ($\dim_{\R} E^* = 2$); choosing a basis $\{\ell_0|_E, \ell_1|_E\}$ defines $\psi_0(u_0) = (u_0, c_0(u_0))$ killing $n = 0, 1$. Because $\varepsilon^2 = 0$, $\mathcal{M}(\psi_0(u_0), (0, v_1), (w_0, 0)) = \RePart(u_0 v_1 w_0)$ is unchanged by $c_0$, so the remaining functionals $\{m_n|_E : n \ge 2\}$ on mode $1$ again span $E^*$, killing two more terms $n = 2, 3$ via $\psi_1(v_0) = (v_0, c_1(v_0))$. Finally, $\mathcal{M}(\psi_0(u_0), \psi_1(v_0), (0, w_1)) = \RePart(u_0 v_0 w_1)$ must be represented by the remaining $R - 4$ terms, and \leanid{single\_complex\_rank\_ge} gives $R - 4 \ge 3$, hence $R \ge 7$.
\end{itemize}

\section{Standard vs.\ Nonstandard Proofs of the Semigroup Step}\label{sec:semigroup}

Every proof of necessity above rests on the continuous one-parameter semigroup theorem: if $\mathcal{A}$ is a Banach algebra over $\R$ and $M : \R \to \mathcal{A}$ satisfies $M(0) = 1$, $M(s + t) = M(s)M(t)$ for all $s, t \ge 0$, and $M$ is continuous on $[0,\infty)$, then there exists $A \in \mathcal{A}$ such that $M(t) = \exp(tA)$ for all $t \ge 0$.

\subsection{Integral smoothing vs.\ Euler's infinitesimal step with homogeneous scale}

\paragraph{Integral smoothing (\leanid{Semigroup.lean}).}
Since $M(0) = 1$ and $M$ is continuous at $0$, there exists $\delta > 0$ such that $I_{\delta} = \int_0^\delta M(u)\,du$ satisfies $\|I_\delta - \delta \cdot 1\| < \delta$, making $I_\delta$ invertible in $\mathcal{A}$ by Neumann series. For $t \ge 0$, changing variables in $M(t) I_\delta = \int_t^{t+\delta} M(u)\,du$ shows by the fundamental theorem of calculus that $M$ is right-differentiable on $[0,\infty)$ with $M'(t) = M(t) A$, where $A = (M(\delta) - 1) I_\delta^{-1}$ commutes with $M(t)$. Differentiating $t \mapsto M(t)\exp(-tA)$ gives $0$, so $M(t) = \exp(tA)$ in \emph{any} Banach algebra.

\paragraph{Finite differences (\leanid{GeneratorBound.lean}, \leanid{FiniteDifferenceSemigroup.lean}).}
Without integration, if $\|M(u) - 1\| \le \frac{1}{2}$ on $[0,\delta]$, geometric-series telescoping bounds the difference quotients $\|(M(\delta/n) - 1)/(\delta/n)\|$. In finite dimensions, compactness extracts a convergent subsequence $(M(h_k) - 1)/h_k \to A$; along $h_k$, the one-sided difference quotients of $M(t)\exp(-tA)$ tend to $0$, and a one-sided mean-value lemma shows $M(t)\exp(-tA) = 1$.

\paragraph{Nonstandard analysis with homogeneous common scale.}
Euler's heuristic sets $A = \operatorname{st}((M(\varepsilon) - 1)/\varepsilon)$ for a positive infinitesimal $\varepsilon > 0$ and iterates $M(\varepsilon)^{\lfloor t/\varepsilon\rfloor} \approx (1 + \varepsilon A)^{\lfloor t/\varepsilon\rfloor} \approx \exp(tA)$. Directly applying $\operatorname{st}$ to $(M(\varepsilon) - 1)/\varepsilon$, however, requires proving beforehand that this quotient is finite (the $124$-line geometric-sum bound above).

All three nonstandard formalizations in this project eliminate that \emph{a priori} bound by passing to \emph{projective / homogeneous coordinates}: put $x = M(\varepsilon) - 1$ (which is infinitesimal by continuity at $0$) and define the common infinitesimal scale
\begin{equation}\label{eq:eta}
  \eta = \varepsilon + \|x\| > 0.
\end{equation}
Because $\frac{\varepsilon}{\eta} + \left\|\frac{x}{\eta}\right\| = 1$ with $\frac{\varepsilon}{\eta} \in [0,1]$ and $\left\|\frac{x}{\eta}\right\| \le 1$, both rescaled quantities are bounded by $1$ \emph{without any analytical estimate}. In a finite-dimensional algebra $\mathcal{A}$, Heine--Borel compactness gives a standard vector $B = \operatorname{st}(x/\eta)$ with $\|B\| \le 1$ and $\varepsilon/\eta \approx a := 1 - \|B\| \in [0,1]$. Now compare $N = \lfloor s / \eta \rfloor$ steps of $M(\varepsilon)$ (arriving at time $N\varepsilon \approx s a$) against $N$ steps of $\exp(\eta B)$ (arriving at time $N\eta \approx s$) for an arbitrary standard $s \ge 0$:
\[
  \frac{\|M(\varepsilon) - \exp(\eta B)\|}{\eta}
  \le \left\|\frac{x}{\eta} - B\right\| + \left\|B - \frac{\exp(\eta B) - 1}{\eta}\right\| \approx 0.
\]
Summing $N$ single-step errors of size $\eta \cdot o(1)$ with $N\eta \le s$ yields $M(s a) = \exp(s B)$ for all standard $s \ge 0$. Could the time coordinate vanish ($a = 0$)? If $a = 0$, then $\exp(s B) = M(0) = 1$ for every $s \ge 0$; differentiating at $s = 0$ gives $B = 0$, contradicting $a + \|B\| = 1$! Hence $a > 0$, and setting $A = B/a$ with $s = t/a$ gives $M(t) = \exp(tA)$.

\subsection{Three nonstandard formalizations and quantitative comparison}

We formalized the homogeneous-scale nonstandard proof in three settings:
\begin{enumerate}[leftmargin=2em,itemsep=2pt]
  \item \textbf{Isabelle/HOL:} \texttt{Isabelle/}\allowbreak\leanid{Semigroup\_NSA.thy} uses \leanid{HOL-Nonstandard\_Analysis}, which provides hyperreals, hypernaturals, star extensions \texttt{*f*}, and the \texttt{transfer} proof method. Single-step error accumulation reduces to a purely algebraic noncommutative power estimate (\leanid{power\_difference\_bound}): if $\|u\|, \|v\| \le e^h$ with $h \ge 0$, then $\|u^n - v^n\| \le n e^{nh}\|u - v\|$, proved by induction from $u^{n+1} - v^{n+1} = u(u^n - v^n) + (u - v)v^n$. Since $\|M(\varepsilon)\| \le 1 + \|x\| \le 1 + \eta \le e^\eta$ and $\|\exp(\eta B)\| \le e^{\eta\|B\|} \le e^\eta$, \texttt{transfer} applies this ordinary inequality directly to a hypernatural step count $N$. (This uses $\|1\| = 1$ from Isabelle's \leanid{real\_normed\_algebra\_1}; in Lean, $\mathcal{L}(\mathcal{W}(\psi))$ has $\|1\| = 0$ when $\psi = 0$, so the Lean theorems avoid assuming $\|1\| = 1$ and telescope steps using compact-interval bounds via \leanid{norm\_replace\_steps\_le} instead.)
  \item \textbf{Lean~4 with ultrafilter sequences:} \texttt{Experiments/}\allowbreak\leanid{NonstandardSemigroup.lean} expresses ``infinitely close'' directly as convergence along $\operatorname{hyperfilter}\ \N$ (since Mathlib defines $\mathrm{Hyperreal} = \operatorname{Germ}(\operatorname{hyperfilter}\ \N)\ \R$ without a transfer tactic), obtaining $B$ in one line via \texttt{IsCompact.}\allowbreak\leanid{ultrafilter\_le\_nhds}.
  \item \textbf{Lean~4 in nonstandard vocabulary:} In \texttt{Experiments/}, \texttt{Ultrapower.lean} and \texttt{Ultrapower\-Semigroup.lean} build a reusable $195$-line ultrapower interface (\leanid{Hyp}, \leanid{std}, \leanid{ext}, $\approx_{\mathrm{h}}$, one-variable transfer lemmas) and write Euler's proof in classical nonstandard notation.
\end{enumerate}

\begin{table}[t]
\centering
\caption{Verified source sizes for the five complete proofs of the continuous one-parameter semigroup step $M(t) = \exp(tA)$. Every project helper lemma in each proof's module dependencies is counted in full (excluding standard Mathlib and Isabelle/HOL library imports).}\label{tab:semigroup}
\smallskip
{\small\setlength{\tabcolsep}{4.5pt}
\begin{tabular}{@{}llrrl@{}}
\toprule
Proof method & Prover / files & Phys. & Nonbl. & Scope \\
\midrule
Integral smoothing (main) & Lean~4 (\texttt{Semigroup.lean}) & 96 & 87 & Any Banach algebra \\
Finite differences & Lean~4 (\texttt{GeneratorBound}, \texttt{FD}) & 256 & 233 & Finite-dim. \\
NSA (ultrafilter seq.) & Lean~4 (\texttt{NonstandardSemigroup}) & 206 & 190 & Finite-dim. \\
NSA (ultrapower API) & Lean~4 (\texttt{Ultrapower*} + shared) & 399 & 344 & Finite-dim. \\
NSA (star + \texttt{transfer}) & Isabelle (\texttt{Semigroup\_NSA.thy}) & 275 & 253 & Finite-dim.\ ($\|1\|=1$) \\
\bottomrule
\end{tabular}}
\end{table}

Table~\ref{tab:semigroup} summarizes the exact line counts. These numbers reflect \emph{library amortization} rather than intrinsic proof complexity: \leanid{Semigroup.lean} ($96$ lines) is shortest and holds in every Banach algebra because Mathlib amortizes thousands of lines of Bochner integration, Neumann series, and the fundamental theorem of calculus behind compact API lemmas, whereas compactness-based proofs are inherently finite-dimensional. Across different proof styles within Lean, the $2\times$ gap between integral smoothing ($96$ lines) and ultrafilter sequences ($206$ lines) remained stable even after eliminating the $124$-line difference-quotient bound via homogeneous coordinates~\eqref{eq:eta}. Moreover, discounting the $195$-line interface \leanid{Ultrapower.lean} as reusable library code leaves $169 + 35 = 204$ lines: without automated transfer, nonstandard vocabulary in Lean merely wraps filter limits. In Isabelle/HOL, \texttt{transfer} genuinely compresses hypernatural induction, but $59$ lines of local library gaps (Heine--Borel vector standard parts and half-line continuity via $M(\max(t,0))$) plus stating transferred formulas at both standard and nonstandard types keep the total at $275$ lines.

\section{Applications to Isotropic and Steerable Filtering and 3-D Propagation}\label{sec:applications}

\subsection{Harmonic ``beach-ball'' differential operators and steerable kernels}

Specializing Proposition~\ref{prop:roots} to equal symbolic factors $z_0 = \dots = z_{d-1} = \partial_x + i\partial_y$ (writing $m$ for the harmonic order and $D_\alpha = \cos\alpha\,\partial_x + \sin\alpha\,\partial_y$ for the directional derivative along angle $\alpha$) gives, for \emph{every} integer $m \ge 1$ (both odd and even),
\begin{equation}\label{eq:beachball}
  \RePart\!\bigl((\partial_x + i\partial_y)^m\bigr)
  = \frac{2^{m-1}}{m} \sum_{k=0}^{m-1} (-1)^k D_{\pi k / m}^m.
\end{equation}
On a regular $2m$-point ring $\theta_n = \pi n / m$ ($n = 0, \dots, 2m-1$) of radius $h \to 0$, antipodal points satisfy $D_{\theta_{k+m}} = -D_{\theta_k}$ while Fourier orthogonality of $(-1)^k = e^{i\pi k}$ annihilates all Taylor orders below $m$. For $f \in C^m$ near $(x, y)$,~\eqref{eq:beachball} recovers---by elementary geometric series rather than Cauchy residues and Chebyshev partial fractions---both the odd-order ($m = 2N+1$) hexagonal/decagonal ``beach-ball'' operators of Niemitalo~\cite{niemitalo2019beachball} and the even-order operators of Perona~\cite{perona1995deformable} in a single formula:
\[
  \lim_{h \to 0} \frac{1}{h^m} \sum_{n=0}^{2m-1} (-1)^n f\!\left(x + h\cos\frac{\pi n}{m},\, y + h\sin\frac{\pi n}{m}\right)
  = \frac{1}{(m-1)!\,2^{m-2}} \RePart\!\bigl((\partial_x + i\partial_y)^m\bigr) f(x,y),
\]
with a half-sector ring rotation $+\frac{\pi}{2m}$ yielding $+\ImPart((\partial_x + i\partial_y)^m)$. Recognizing $\partial_{\bar z} = \frac{1}{2}(\partial_x + i\partial_y)$ as the Cauchy--Riemann operator also explains why applying $\partial_{\bar z}^m$ to \emph{any} smooth radial kernel $\varphi(x^2 + y^2) = \varphi(z\bar z)$ produces no lower-order radial derivatives: because $\partial_{\bar z}(z) = 0$, induction gives
\begin{equation}\label{eq:wirtinger_radial}
  \left(\frac{\partial_x + i\partial_y}{2}\right)^m \varphi(x^2 + y^2)
  = (x + iy)^m\,\varphi^{(m)}(x^2 + y^2)
  = r^m e^{i m \theta}\,\varphi^{(m)}(r^2).
\end{equation}

When $\varphi(s)$ is a well-formed multi-Gaussian profile with $r$ real and $c$ nonreal upper-half-plane rates (all $\lambda_i \ne 0$, approximating, e.g., an ideal circular low-pass filter~\cite{niemitalo2019lowpass}), $\varphi^{(m)}(s)$ is \emph{again} a well-formed multi-Gaussian profile at the same rates (weighted by $\lambda_i^m$) of continuous 2-way rank $R = r + 2c$. Analytically differentiating the continuous-space kernel on $\R^2$---avoiding discrete finite-difference distortion up to Nyquist~\cite{niemitalo2019angle}---absorbs the monomials of $\RePart((x+iy)^m)$ and $\ImPart((x+iy)^m)$ into the 1-D Gaussian factors, giving exact \emph{continuous-domain} 2-way real separable ranks on $\R^2$:
\[
  \bigl(\lfloor m/2 \rfloor + 1\bigr)(r + 2c) \quad \text{for } \cos(m\theta), \qquad
  \lceil m/2 \rceil (r + 2c) \quad \text{for } \sin(m\theta).
\]
Sampling those 1-D factors on a finite grid yields a discrete separable realization with at most that many terms (grid aliasing at special rates such as $\ImPart\lambda_i \in \pi\mathbb{Z}$ can collapse the sampled discrete rank further).

\paragraph{Continuous-domain pass-count bounds for simultaneous steerable pairs.}
In applications such as orientation estimation~\cite{niemitalo2019angle}, one evaluates both continuous-domain convolutions $(C_m * I, S_m * I)$ of a harmonic pair simultaneously. Can sharing intermediate 1-D continuous convolutions reduce the pass count below separate evaluation $2(\lfloor m/2\rfloor + 1)R + 2\lceil m/2\rceil R = 2(m + 1)R$? No: if any two-stage 1-D schedule on $\R^2$ splits $(C_m, S_m) = (A_c + B_c, A_s + B_s)$ into $P_y$ shared vertical passes followed by $\operatorname{rank}(A_c) + \operatorname{rank}(A_s)$ horizontal passes plus $P_x$ shared horizontal passes followed by $\operatorname{rank}(B_c) + \operatorname{rank}(B_s)$ vertical passes, then
\begin{align*}
  \dim\operatorname{rowSpan}_x(C_m, S_m) &\le \operatorname{rank}(A_c) + \operatorname{rank}(A_s) + P_x,\\
  \dim\operatorname{colSpan}_y(C_m, S_m) &\le P_y + \operatorname{rank}(B_c) + \operatorname{rank}(B_s).
\end{align*}
Because $C_m$ and $S_m$ on $\R^2$ span all powers $\{x^0,\dots,x^m\}g_n(x)$ in $x$ and $\{y^0,\dots,y^m\}g_n(y)$ in $y$ (for $m = 1$, $h_x$ is odd-in-$x$/even-in-$y$ while $h_y$ is even-in-$x$/odd-in-$y$), adding the two inequalities proves that \emph{every} two-stage continuous 1-D schedule requires at least $(m+1)R + (m+1)R = 2(m+1)R$ passes ($4R$ passes for $(h_x, h_y)$). Savings instead arise in three ways: (i)~definite even/odd parity of every 1-D factor $u^k g_n(u)$ halves the multiplications per sampled pass via $I[y+\ell] \pm I[y-\ell]$ folding while fusing input loads; (ii)~across \emph{multiple} harmonic orders $m = 0,\dots,M$, vertical passes $\{I *_y (y^k g_n)\}_{k=0}^M$ are shared linearly ($5R$ passes instead of $6R$ for low-pass plus gradient; $10R$ instead of $12R$ for $m = 1$ plus $m = 3$); and (iii)~for the $d$-dimensional gradient $(h_{x_0},\dots,h_{x_{d-1}})$ of a multi-Gaussian radial kernel of $d$-way rank $R_d = r + d c$ ($d \ge 3$), recursively splitting $\{0,\dots,d-1\}$ into halves of sizes $\lfloor d/2\rfloor$ and $\lceil d/2\rceil$ and pre-convolving each half with the undifferentiated factors reduces $d^2 R_d$ passes to $T(d) R_d$ with $T(1) = 1$ and $T(d) = T(\lfloor d/2\rfloor) + T(\lceil d/2\rceil) + d \le d(\lceil \log_2 d \rceil + 1)$ ($8R_3$ vs.\ $9R_3$ in $d = 3$; $12R_4$ vs.\ $16R_4$ in $d = 4$).

\subsection{Three-dimensional filtering, wave propagation, and concrete operation counts}

In three dimensions ($d = 3$), direct spatial convolution with an $L \times L \times L$ isotropic kernel---such as a cryo-EM/tomography bandpass filter or a coarse-grid dispersion-free time-step propagator for the band-limited constant-speed acoustic wave equation, whose spatial-frequency symbol depends radially on $\|\xi\|^2$---scales as $L^3$ per voxel. Approximating the radial profile by $r$ real and $c$ complex-conjugate Gaussians yields exact real 3-way separable rank $r + 3c$ by Theorem~\ref{thm:real_nd}: Proposition~\ref{prop:roots} realizes each complex Gaussian in $3$ real 3-way terms ($9$ real 1-D passes on real-valued volume buffers, vs.\ $2^{3-1} = 4$ terms / $12$ passes from naive real-and-imaginary expansion), and polynomial-oscillatory shells ($\deg P = 1$) take $7$ rather than $9$ real 3-way terms via Proposition~\ref{prop:poly_roots}.

\begin{table}[t]
\centering
\caption{Comparison of exact continuous-domain real separable rank $R$, 1-D passes, and scalar multiply--adds per output sample ($d R L$) against direct spatial convolution ($L^d$) and naive real-part expansion at tap widths $L = 41$ and $L = 161$ (from~\cite{niemitalo2019angle}).}\label{tab:examples}
\smallskip
{\footnotesize\setlength{\tabcolsep}{3.2pt}
\begin{tabular}{@{}llrrrrr@{}}
\toprule
Kernel family & Params & Direct ($41 / 161$) & Naive & Exact $R$ & Passes & Ops ($41 / 161$) \\
\midrule
2-D circular low-pass $W$ & $d{=}2,\,c{=}2$ & $1{,}681$ / $25{,}921$ & $8$ & $4$ & $8$ & $328$ / $1{,}288$ \\
2-D gradient $(h_x,h_y)$ & $d{=}2,\,m{=}1,\,c{=}2$ & $3{,}362$ / $51{,}842$ & $16$ & $4{+}4$ & $16$ & $656$ / $2{,}576$ \\
2-D beach ball $(C_3,S_3)$ & $d{=}2,\,m{=}3,\,c{=}2$ & $3{,}362$ / $51{,}842$ & $32$ & $8{+}8$ & $32$ ($40^*$) & $1{,}312$ / $5{,}152$ \\
3-D Gaussian prop. & $d{=}3,\,r{=}1,\,c{=}2$ & $68{,}921$ / $4.17{\times}10^6$ & $27$ & $7$ & $21$ & $861$ / $3{,}381$ \\
3-D osc.\ shell ($\deg{=}1$) & $d{=}3,\,\deg{=}1,\,c{=}1$ & $68{,}921$ / $4.17{\times}10^6$ & $27$ & $7$ & $21$ & $861$ / $3{,}381$ \\
\bottomrule
\end{tabular}}
\end{table}

Table~\ref{tab:examples} illustrates concrete operation counts (before an additional $2\times$ reduction in scalar multiplications from even/odd tap-parity folding; $40^*$ denotes the combined pass count for $m = 1$ plus $m = 3$, vs.\ $48$ unshared). On the $N = 161$ ($\omega_c = \pi/64$) orientation-estimation benchmark of~\cite{niemitalo2019angle}, replacing direct $2 \times 161^2$ spatial convolution by a two-complex-Gaussian analytical gradient pair reduces arithmetic by $20\times$ ($40\times$ with tap-parity folding) without full-image 2-D FFT buffers; in 3-D at $L = 41$, the $r + 3c = 7$ decomposition reduces $68{,}921$ taps/voxel to $861$ ($80\times$), saving $22\%$ over naive $1 + 4c = 9$ real/imaginary expansion.

\section{Formalization Architecture and Verification Audit}\label{sec:formalization}

\begin{table}[t]
\centering
\caption{Structure and source line counts of the \texttt{KernelRankLean} library.}\label{tab:modules}
\smallskip
{\small\setlength{\tabcolsep}{3.5pt}
\begin{tabular}{@{}llrr@{}}
\toprule
Module(s) & Mathematical content & Phys. & Nonbl. \\
\midrule
\texttt{Defs}, \texttt{KernelRankLean} & Multi-way rank, surjections, slices, costs & 213 & 180 \\
\texttt{RowSpan}, \texttt{Complexification} & 2-way row-span rank; $\R \to \C$ invariance & 221 & 189 \\
\texttt{Translations} & $d$-D radial-to-Hankel reduction; \texttt{translateSpan} & 92 & 80 \\
\texttt{ExpPoly}, \texttt{RealExpPoly} & Exponential-polynomial uniqueness & 198 & 170 \\
\texttt{Semigroup}, \texttt{TranslationSemigroup} & Banach semigroups; shift semigroup & 262 & 219 \\
\texttt{Necessity}, \texttt{Representation} & Generalized eigenvectors; conjugate pairing & 222 & 198 \\
\texttt{Differential}, \texttt{ComplexRank} & Half-line operators $q(D)$; 2-D complex rank & 318 & 284 \\
\texttt{Annihilator} & Minimal constant-coefficient annihilator & 69 & 62 \\
\texttt{MultiComplexRank} & Exact $d$-way complex rank $\sum_i ((d-1)\deg P_i + 1)$ & 1,015 & 978 \\
\texttt{MultiRealRank} & Exact $d$-way real rank (\texttt{multiTotalCost}) & 2,562 & 2,521 \\
\texttt{Main} & Unified $d$-D and 2-D main theorems & 152 & 140 \\
\midrule
\textbf{Total (\texttt{KernelRankLean})} & & \textbf{5,324} & \textbf{5,021} \\
\bottomrule
\end{tabular}}
\end{table}

Table~\ref{tab:modules} breaks down the \texttt{KernelRankLean} library ($5{,}324$ physical lines, $5{,}021$ nonblank lines; earlier stages comprised $3{,}822$ physical lines at commit \texttt{6aabb4d} for $d$-D complex and multi-Gaussian ranks, $1{,}587$ lines at commit \texttt{cae13d6} for 2-D plus differential annihilators, and $1{,}340$ lines at commit \texttt{240dea0} for 2-D classification alone). Three formalization choices proved decisive in Lean~4:

\begin{enumerate}[leftmargin=2em,itemsep=2pt]
  \item \textbf{Dependent multi-way definitions with homogeneous two-coordinate specialization:} Defining \leanid{IsSumOfProducts} and \leanid{separableRank} over a general dependent family $\prod_{i \in \iota} X_i$ unifies 2-D and $d$-D statements, whereas specializing \leanid{rowSpan} in \leanid{RowSpan.lean} to homogeneous $\operatorname{Fin} 2 \to \alpha$ (with a pattern-matched constructor \leanid{fin2Vec}) avoids expensive higher-order unification timeouts when rewriting on $\operatorname{Fin} 2 \to \Rnn$.
  \item \textbf{Explicit uniformity selection on \leanid{translateSpan}:} Transferring a norm from $\C^n$ onto the finite-dimensional subspace $\mathcal{W}(\psi) \le (\Rnn \to \C)$ creates a second \leanid{UniformSpace} instance alongside the subtype uniformity; explicitly installing the norm-induced uniformity as the local instance prevents typeclass diamond timeouts while letting \texttt{LinearMap.}\allowbreak\leanid{continuous\_of\_finiteDimensional} supply automatic continuity.
  \item \textbf{Cayley--Hamilton divisibility in place of cyclic vectors:} Bounding $\dim_{\C}\mathcal{W}(\psi)$ from below via divisibility of $\chi_A(X)$ by $\prod_\mu (X - \mu)^{\deg P_\mu + 1}$ eliminated both coefficient-synthesis operators and cyclic-vector constructions, shrinking the 2-D proof plus differential-annihilator extension by $93$ lines while yielding Theorem~\ref{thm:annihilator} directly.
\end{enumerate}

\paragraph{Verification audit.}
A single script (\texttt{scripts/check}, also wired as the repository's pre-commit hook) builds both \texttt{KernelRankLean} and \texttt{Experiments} via \texttt{lake build} (which includes every \texttt{.lean} module in the repository), checks the compiler log to ensure no declaration uses \texttt{sorry}, runs \texttt{\#print axioms} on $17$ top-level theorems to confirm dependence solely on Lean's three standard classical axioms (\texttt{propext}, \texttt{Classical.choice}, \texttt{Quot.sound}), builds the Isabelle2026 session \leanid{KernelRank\_NSA}, and audits \texttt{Semigroup\_NSA.}\allowbreak\leanid{semigroup\_eq\_exp} via \leanid{Thm\_Deps.all\_oracles} (catching \leanid{skip\_proof}/\texttt{sorry} even under \leanid{quick\_and\_dirty}) while rejecting any \texttt{axiomatization} or \texttt{axioms} declarations.

\section*{Acknowledgments}

This work builds directly on problems, mathematical formulations, and software developed by Olli Niemitalo, including \texttt{kernel-optimizer}~\cite{niemitalo_kernel_optimizer} for complex-Gaussian approximations of 2-D circularly symmetric blur kernels and his investigations of isotropic low-pass filters, steerable gradient orientation estimation, and planar harmonic differential operators~\cite{niemitalo2019angle,niemitalo2019beachball,niemitalo2019lowpass}. Olli Niemitalo initiated and guided the \texttt{kernel-rank-lean} repository, proposed comparing standard and Isabelle/HOL nonstandard proofs of the semigroup step, and suggested extending the formalization to $d$ dimensions. We also thank the aiXiv Official Agent reviewer on version~1.0, whose question regarding higher-degree real polynomial profiles in $d \ge 3$ directly motivated Theorem~\ref{thm:real_nd} and Proposition~\ref{prop:poly_roots}. Both the repository and this manuscript are released under CC0~1.0 Universal.

\begin{thebibliography}{99}

\bibitem{alder1981alg}
A.~Alder and V.~Strassen,
``On the algorithmic complexity of associative algebras,''
\emph{Theoretical Computer Science} \textbf{15}(2), 201--211 (1981).

\bibitem{beylkin2005exponentials}
G.~Beylkin and L.~Monz\'on,
``On approximation of functions by exponential sums,''
\emph{Applied and Computational Harmonic Analysis} \textbf{19}(1), 17--48 (2005).

\bibitem{comon1996decomposition}
P.~Comon and B.~Mourrain,
``Decomposition of quantics in sums of powers of linear forms,''
\emph{Signal Processing} \textbf{53}(2--3), 93--107 (1996).

\bibitem{fleuriot2000mechanizing}
J.~D.~Fleuriot,
``On the mechanization of real analysis in Isabelle/HOL,''
in \emph{Theorem Proving in Higher Order Logics (TPHOLs 2000)}, LNCS \textbf{1869}, pp.~145--161 (2000).

\bibitem{freeman1991steerable}
W.~T.~Freeman and E.~H.~Adelson,
``The design and use of steerable filters,''
\emph{IEEE Transactions on Pattern Analysis and Machine Intelligence} \textbf{13}(9), 891--906 (1991).

\bibitem{holzl2013type}
J.~H\"olzl, F.~Immler, and B.~Huffman,
``Type classes and filters for mathematical analysis in Isabelle/HOL,''
in \emph{Interactive Theorem Proving (ITP 2013)}, LNCS \textbf{7998}, pp.~279--294 (2013).

\bibitem{koenderink1992generic}
J.~J.~Koenderink and A.~J.~van~Doorn,
``Generic neighborhood operators,''
\emph{IEEE Transactions on Pattern Analysis and Machine Intelligence} \textbf{14}(6), 597--605 (1992).

\bibitem{lampert2006anisotropic}
C.~H.~Lampert and O.~Wirjadi,
``An optimal nonorthogonal separation of the anisotropic Gaussian convolution filter,''
\emph{IEEE Transactions on Image Processing} \textbf{15}(11), 3501--3513 (2006).

\bibitem{landsberg2012tensors}
J.~M.~Landsberg,
\emph{Tensors: Geometry and Applications},
Graduate Studies in Mathematics \textbf{128}, American Mathematical Society (2012).

\bibitem{mathlib2020}
The mathlib Community,
``The Lean mathematical library,''
in \emph{Proceedings of the 9th ACM SIGPLAN International Conference on Certified Programs and Proofs (CPP 2020)}, pp.~367--381 (2020).

\bibitem{niemitalo2019angle}
O.~Niemitalo,
``Auto detection of rotation angle on arbitrary image with orthogonal features'' (answer),
\emph{Signal Processing Stack Exchange} (2019),
\url{https://dsp.stackexchange.com/a/58287/15347}.

\bibitem{niemitalo2019beachball}
O.~Niemitalo,
``Beach-ball like differential operators of a two-dimensional function'' (question; with answer by K.~B.~Dave),
\emph{Mathematics Stack Exchange} (2019),
\url{https://math.stackexchange.com/q/3232122}.

\bibitem{niemitalo2019lowpass}
O.~Niemitalo,
``Efficient implementation of 2-D circularly symmetric low-pass filter,''
\emph{Signal Processing Stack Exchange} (2019),
\url{https://dsp.stackexchange.com/q/58449}.

\bibitem{niemitalo_kernel_optimizer}
O.~Niemitalo,
\emph{kernel-optimizer: Circularly symmetric 2-D blur kernel approximation by complex Gaussians},
GitHub repository, \url{https://github.com/OlliNiemitalo/kernel-optimizer}.

\bibitem{perona1995deformable}
P.~Perona,
``Deformable kernels for early vision,''
\emph{IEEE Transactions on Pattern Analysis and Machine Intelligence} \textbf{17}(5), 488--499 (1995).

\bibitem{qi2016semialgebraic}
Y.~Qi, P.~Comon, and L.-H.~Lim,
``Semialgebraic geometry of nonnegative tensor rank,''
\emph{SIAM Journal on Matrix Analysis and Applications} \textbf{37}(4), 1556--1580 (2016).

\bibitem{reznick2013length}
B.~Reznick,
``On the length of binary forms,''
in \emph{ Quadratic and Higher Degree Forms}, Developments in Mathematics \textbf{31}, pp.~207--232 (2013).

\bibitem{schoenberg1938metric}
I.~J.~Schoenberg,
``Metric spaces and completely monotone functions,''
\emph{Annals of Mathematics} \textbf{39}(4), 811--841 (1938).

\bibitem{unser2013steerable}
M.~Unser and N.~Chenouard,
``A unifying parametric framework for 2D steerable wavelet transforms,''
\emph{SIAM Journal on Imaging Sciences} \textbf{6}(1), 102--135 (2013).

\end{thebibliography}

\end{document}
